% GENERATED FILE. Edit canonical lessons, then run npm run generate:books.
% canonical-source-hash: fnv1a64:e9e9703d76e01b6d
% canonical-lessons: BN-C181-nihshvas-neoya, BN-C181-kapa, BN-C181-bheja, BN-C181-mukhostho-kora, BN-C181-bojhano

\chapter{To breathe, To shiver, To get wet, To memorise, To explain}
\label{ch:bn-to-breathe-to-shiver-to-get-wet-to-memorise-to-explain}
\hlchaptermodality{181}

\begin{chapteropening}
\textbf{By the end of this chapter:} \emph{I can say to breathe, to shiver, to get wet, to memorise and to explain.}

Complete the last lesson of chapter 181: I can say to breathe, to shiver, to get wet, to memorise and to explain.
\end{chapteropening}

\section[\texorpdfstring{\bn{নিঃশ্বাস} \bn{নেওয়া} (niḥshvās neoyā)}{নিঃশ্বাস নেওয়া (niḥshvās neoyā)}]{\bn{নিঃশ্বাস} \bn{নেওয়া} (niḥshvās neoyā) — to breathe}

\label{lesson:BN-C181-nihshvas-neoya}



\begin{quote}
Before the new one: say the Bengali for to water, then the Bengali for to print.
\end{quote}

\subsection*{You'll want to know: \bn{নিঃশ্বাস} \bn{নেওয়া}}
\textbf{\bn{নিঃশ্বাস} \bn{নেওয়া}} — \emph{niḥshvās neoyā} — \textquotedblleft{}to breathe\textquotedblright{}.

\begin{grammarlens}[title={What you've built}]
One more word for this chapter.
\end{grammarlens}

\subsection*{Your turn}
\begin{itemize}
\raggedright
  \item \emph{Say it:} \emph{niḥshvās neoyā}
  \item \emph{Say it:} \emph{niḥshvās neoyā}, once more
  \item \emph{Say it:} say \emph{niḥshvās neoyā}
  \item \emph{Recall:} say the Bengali for to water, then the Bengali for to print, then say \emph{niḥshvās neoyā} again
\end{itemize}

\begin{tcolorbox}[breakable,colback=teal!4,colframe=teal!35!black,title={Before you move on}]
What does \bn{নিঃশ্বাস} \bn{নেওয়া} mean? (\textquotedblleft{}To breathe\textquotedblright{}.) Say it once more.
\end{tcolorbox}
\section[\texorpdfstring{\bn{কাঁপা} (kā̃pā)}{কাঁপা (kā̃pā)}]{\bn{কাঁপা} (kā̃pā) — to shiver}

\label{lesson:BN-C181-kapa}



\begin{quote}
Before the new one: say the Bengali for to print, then the Bengali for to breathe.
\end{quote}

\subsection*{You'll want to know: \bn{কাঁপা}}
\textbf{\bn{কাঁপা}} — \emph{kā̃pā} — \textquotedblleft{}to shiver\textquotedblright{}.

\begin{grammarlens}[title={What you've built}]
One more word for this chapter.
\end{grammarlens}

\subsection*{Your turn}
\begin{itemize}
\raggedright
  \item \emph{Say it:} \emph{kā̃pā}
  \item \emph{Say it:} \emph{kā̃pā}, once more
  \item \emph{Say it:} say \emph{kā̃pā}
  \item \emph{Recall:} say the Bengali for to print, then the Bengali for to breathe, then say \emph{kā̃pā} again
\end{itemize}

\begin{tcolorbox}[breakable,colback=teal!4,colframe=teal!35!black,title={Before you move on}]
What does \bn{কাঁপা} mean? (\textquotedblleft{}To shiver\textquotedblright{}.) Say it once more.
\end{tcolorbox}
\section[\texorpdfstring{\bn{ভেজা} (bhejā)}{ভেজা (bhejā)}]{\bn{ভেজা} (bhejā) — to get wet}

\label{lesson:BN-C181-bheja}



\begin{quote}
Before the new one: say the Bengali for to breathe, then the Bengali for to shiver.
\end{quote}

\subsection*{You'll want to know: \bn{ভেজা}}
\textbf{\bn{ভেজা}} — \emph{bhejā} — \textquotedblleft{}to get wet\textquotedblright{}.

\begin{grammarlens}[title={What you've built}]
One more word for this chapter.
\end{grammarlens}

\subsection*{Your turn}
\begin{itemize}
\raggedright
  \item \emph{Say it:} \emph{bhejā}
  \item \emph{Say it:} \emph{bhejā}, once more
  \item \emph{Say it:} say \emph{bhejā}
  \item \emph{Recall:} say the Bengali for to breathe, then the Bengali for to shiver, then say \emph{bhejā} again
\end{itemize}

\begin{tcolorbox}[breakable,colback=teal!4,colframe=teal!35!black,title={Before you move on}]
What does \bn{ভেজা} mean? (\textquotedblleft{}To get wet\textquotedblright{}.) Say it once more.
\end{tcolorbox}
\section[\texorpdfstring{\bn{মুখস্থ} \bn{করা} (mukhôsthô kôrā)}{মুখস্থ করা (mukhôsthô kôrā)}]{\bn{মুখস্থ} \bn{করা} (mukhôsthô kôrā) — to memorise}

\label{lesson:BN-C181-mukhostho-kora}



\begin{quote}
Before the new one: say the Bengali for to shiver, then the Bengali for to get wet.
\end{quote}

\subsection*{You'll want to know: \bn{মুখস্থ} \bn{করা}}
\textbf{\bn{মুখস্থ} \bn{করা}} — \emph{mukhôsthô kôrā} — \textquotedblleft{}to memorise\textquotedblright{}.

\begin{grammarlens}[title={What you've built}]
One more word for this chapter.
\end{grammarlens}

\subsection*{Your turn}
\begin{itemize}
\raggedright
  \item \emph{Say it:} \emph{mukhôsthô kôrā}
  \item \emph{Say it:} \emph{mukhôsthô kôrā}, once more
  \item \emph{Say it:} say \emph{mukhôsthô kôrā}
  \item \emph{Recall:} say the Bengali for to shiver, then the Bengali for to get wet, then say \emph{mukhôsthô kôrā} again
\end{itemize}

\begin{tcolorbox}[breakable,colback=teal!4,colframe=teal!35!black,title={Before you move on}]
What does \bn{মুখস্থ} \bn{করা} mean? (\textquotedblleft{}To memorise\textquotedblright{}.) Say it once more.
\end{tcolorbox}
\section[\texorpdfstring{\bn{বোঝানো} (bojhāno)}{বোঝানো (bojhāno)}]{\bn{বোঝানো} (bojhāno) — to explain}

\label{lesson:BN-C181-bojhano}



\begin{quote}
Before the new one: say the Bengali for to get wet, then the Bengali for to memorise.
\end{quote}

\subsection*{You'll want to know: \bn{বোঝানো}}
\textbf{\bn{বোঝানো}} — \emph{bojhāno} — \textquotedblleft{}to explain\textquotedblright{}.

\begin{grammarlens}[title={What you've built}]
One more word for this chapter.
\end{grammarlens}

\subsection*{Your turn}
\begin{itemize}
\raggedright
  \item \emph{Say it:} \emph{bojhāno}
  \item \emph{Say it:} \emph{bojhāno}, once more
  \item \emph{Say it:} say \emph{bojhāno}
  \item \emph{Recall:} say the Bengali for to get wet, then the Bengali for to memorise, then say \emph{bojhāno} again
\end{itemize}

\begin{tcolorbox}[breakable,colback=teal!4,colframe=teal!35!black,title={Before you move on}]
What does \bn{বোঝানো} mean? (\textquotedblleft{}To explain\textquotedblright{}.) Say it once more.
\end{tcolorbox}
